\begin{theorem}[Jensen’s inequality]
Let \(f:\mathbb{R}\to\mathbb{R}\) be convex, i.e.
\[
\forall x,y\in\mathbb{R},\;\forall t\in[0,1]\qquad
f\bigl(tx+(1-t)y\bigr)\le t\,f(x)+(1-t)f(y). \tag{C}
\]
Let \(x_{1},\dots,x_{n}\in\mathbb{R}\) and weights \(\lambda_{1},\dots,\lambda_{n}\) satisfy
\(\lambda_i\ge0\) and \(\sum_{i=1}^{n}\lambda_i=1\).
Then
\[
f\!\Bigl(\sum_{i=1}^{n}\lambda_i x_i\Bigr)
   \le \sum_{i=1}^{n}\lambda_i f(x_i). \tag{J}
\]
\end{theorem}

\begin{proof}
\begin{enumerate}
\item[\textbf{1.}] \textbf{Weight conditions.} The hypotheses imply that
\[
X:=\sum_{i=1}^{n}\lambda_i x_i
\]
is a convex combination of the numbers \(x_i\).

\item[\textbf{2.}] \textbf{Base case \(n=2\).}  For two points write
\(\lambda_{1}=t,\ \lambda_{2}=1-t\) with \(t\in[0,1]\).  By (C),
\[
f(t x_{1}+(1-t)x_{2})\le t\,f(x_{1})+(1-t)f(x_{2}),
\]
which is exactly (J) for \(n=2\).

\item[\textbf{3.}] \textbf{Induction hypothesis.}  Assume that for a fixed
\(k\ge2\) the statement (J) holds for every collection of \(k\) points and
any non‑negative weights summing to one; i.e. for any \(y_{1},\dots ,y_{k}\)
and any \(\mu_{1},\dots ,\mu_{k}\) with \(\mu_i\ge0\) and \(\sum_{i=1}^{k}\mu_i=1\),
\[
f\!\Bigl(\sum_{i=1}^{k}\mu_i y_i\Bigr)
   \le \sum_{i=1}^{k}\mu_i f(y_i). \tag{IH}
\]

\item[\textbf{4.}] \textbf{Prepare the \((k+1)\)-point case.}  Let
\(x_{1},\dots ,x_{k+1}\) and weights \(\lambda_{1},\dots ,\lambda_{k+1}\) satisfy the
same conditions.  Set
\[
\alpha:=\lambda_{k+1},\qquad
\beta:=1-\alpha=\sum_{i=1}^{k}\lambda_i .
\]
Because the \(\lambda_i\) are non‑negative, \(\beta\in[0,1]\).  If \(\beta=0\) then
\(\lambda_{k+1}=1\) and the desired inequality reduces to the trivial equality
\(f(x_{k+1})=f(x_{k+1})\).  Otherwise define
\[
y:=\frac{1}{\beta}\sum_{i=1}^{k}\lambda_i x_i . \tag{4}
\]
The coefficients \(\lambda_i/\beta\) are non‑negative and sum to one, so (4) is a convex
combination of the first \(k\) points.

\item[\textbf{5.}] \textbf{Apply the induction hypothesis to the first \(k\) points.}
Using (IH) with the points \(x_{1},\dots ,x_{k}\) and the weights
\(\lambda_i/\beta\) we obtain
\[
f(y)=f\!\Bigl(\sum_{i=1}^{k}\frac{\lambda_i}{\beta}x_i\Bigr)
   \le \sum_{i=1}^{k}\frac{\lambda_i}{\beta}\,f(x_i). \tag{5}
\]

\item[\textbf{6.}] \textbf{Rewrite the whole convex combination as a two‑point
convex combination.}  Directly from the definitions,
\begin{align*}
X&:=\sum_{i=1}^{k+1}\lambda_i x_i \\
  &=\beta\Bigl(\frac{1}{\beta}\sum_{i=1}^{k}\lambda_i x_i\Bigr)+\alpha x_{k+1}\\
  &=\beta y+\alpha x_{k+1}. \tag{6}
\end{align*}
Here \(\beta,\alpha\in[0,1]\) and \(\beta+\alpha=1\).

\item[\textbf{7.}] \textbf{Apply convexity to the pair \((y,x_{k+1})\).}
Using (C) with \(t=\beta\) gives
\[
f(X)=f(\beta y+\alpha x_{k+1})
   \le \beta\,f(y)+\alpha\,f(x_{k+1}). \tag{7}
\]

\item[\textbf{8.}] \textbf{Insert the bound (5) for \(f(y)\) into (7).}
Combining (5) and (7),
\begin{align*}
f(X) &\le \beta\Bigl(\sum_{i=1}^{k}\frac{\lambda_i}{\beta}f(x_i)\Bigr)
        +\alpha f(x_{k+1})\\
     &= \sum_{i=1}^{k}\lambda_i f(x_i)+\lambda_{k+1}f(x_{k+1})\\
     &= \sum_{i=1}^{k+1}\lambda_i f(x_i). \tag{8}
\end{align*}
This is precisely the desired inequality for \(k+1\) points.

\item[\textbf{9.}] \textbf{Induction conclusion.}  The base case \(n=2\) holds,
and the induction step shows that truth for \(k\) implies truth for \(k+1\).
Hence, by mathematical induction, (J) is valid for every integer
\(n\ge2\).  The case \(n=1\) is immediate because both sides equal
\(f(x_{1})\).

\item[\textbf{10.}] \textbf{Equality condition.}  Equality in (C) occurs exactly
when either the weight is \(0\) or \(1\) (so one point is not used) or when
\(f\) is affine on the segment joining the two points.  Propagating this
through the induction argument yields equality in Jensen’s inequality
precisely when all points with positive weight are equal, or when \(f\) is
affine on the convex hull of \(\{x_{1},\dots ,x_{n}\}\).

\end{enumerate}
Thus, for any convex function \(f\) and any convex combination of real
numbers, Jensen’s inequality holds.
\end{proof}